\documentclass[11pt]{article}
\usepackage[margin=1in]{geometry}
\usepackage{amsmath,amssymb,amsthm,mathtools}
\usepackage[hidelinks]{hyperref}
\newtheorem{theorem}{Theorem}
\newtheorem{lemma}[theorem]{Lemma}
\newtheorem{corollary}[theorem]{Corollary}
\title{An effective saddle annulus for the exponent-six factorial polynomials}
\author{Internal analytic derivation; no priority claim}
\date{17 November 2026}
\begin{document}
\maketitle
Put
\[
 P_n(z)=\sum_{k=0}^n\binom nk(6k)!z^{n-k},\qquad
 I_n(w)=\int_0^\infty(w+s^6)^n e^{-ns}\,ds.
\]
Then $P_n(n^6w)=n^{6n+1}I_n(w)$. The overarching objective remains a
proof or counterexample for Factorial and/or Strong Factorial. The focused
exponent-six endpoint is effective eventual simplicity, followed by an exact
finite-prefix check. This note closes two geometric regions in that analytic
programme; it does not prove simplicity, give its cutoff, or verify a prefix.
No short-interval or density theorem is used.
The subsequent \href{EFFECTIVE_EXPONENT_SIX_SIMPLICITY.tex}{effective simplicity proof}
now completes the ordinary exponent-six two-term family at every index,
with simplicity cutoff $10000$ and an exact finite prefix. The annulus
results below remain valid as the preceding localization step.

\section{The effective conclusion}
\begin{theorem}\label{thm:annulus}
For every $n\ge1$, every zero of $P_n(n^6w)$ satisfies $|w|<2800$.
For every $n\ge10^6$, the stronger enclosure is
\[
                         1<|w|<2800.                   \tag{1}
\]
More precisely, for $|w|\le1$, let $s_0=s_0(w)$ be the unique solution of
$w=s^5(6-s)$ with $|s_0-6|<10^{-3}$. Put
\[
 \Phi_0=\log(w+s_0^6)-s_0,\qquad a=1-5/s_0,
\]
using the analytic logarithmic chart near $s=6$. For $n\ge10^6$,
\[
 \left|I_n(w)-e^{n\Phi_0}\sqrt{\frac{2\pi}{na}}\right|
 <\frac{e^{n\operatorname{Re}\Phi_0}}{\sqrt n},\qquad
 |I_n(w)|>\frac{4e^{n\operatorname{Re}\Phi_0}}{\sqrt n}. \tag{2}
\]
The square root is the branch obtained from $a=1/6$ at $w=0$.
The integer $10^6$ is a cutoff for this inner zero exclusion only.
\end{theorem}

\section{The outer enclosure}
Write $a_k=\binom nk(6k)!$. Monotone integral bounds on $\log(k!)$ give
$k!\ge(k/e)^k$ and $(6k)!\le6ek(6k/e)^{6k}$. For $3\le k\le n$,
\[
 \frac{a_k^{1/k}}{n^6}
 \le 6^6e^{-5}(k/n)^5(6ek)^{1/k}
 <4\cdot6^6(3/8)^5=\frac{177147}{128}<1400.
\]
Here $(6et)^{1/t}$ decreases for $t\ge3$, and
$(18e)^{1/3}<4$ and $8/3<e<3$. The other cases satisfy
$a_1/n^6\le720$ and, when $n\ge2$,
$a_2^{1/2}/n^6\le\sqrt{12!/2}/32<500$.
Thus, for $|z|\ge2800n^6$,
\[
 \left|\sum_{k=1}^n a_kz^{-k}\right|<\sum_{k=1}^n2^{-k}<1.
\]
This proves the outer assertion without an asymptotic argument.

\section{The inner disk: the saddle and a local Gaussian estimate}
Set $\delta=1/1000$, $r=1/4$ and $G(x)=6\log x-x$.
On $|s-6|=\delta$,
$|s^5(6-s)|\ge(6-\delta)^5\delta>1$.
Rouch\'e's theorem gives the unique $s_0=\alpha+i\beta$ in that disk,
for every $|w|\le1$. The saddle is simple and depends analytically on $w$.
Choose near it
\[
 \Phi(s,w)=6\log s+\log(1+w/s^6)-s.
\]
Then $\Phi_s(s_0,w)=0$, $\Phi_{ss}(s_0,w)=-a$, and
\[
 \operatorname{Re}a>1/7,\qquad |a|<1/5.                 \tag{3}
\]
The first bound follows from
$\operatorname{Re}a\ge1-5/(6-\delta)>1/7$; the second follows from
$|a|\le(1+\delta)/(6-\delta)<1/5$.
Writing $S=\operatorname{Re}\Phi_0$ and using $w+s_0^6=6s_0^5$ gives
\[
 S\ge G(6)+5\log(1-\delta/6)-\delta
    >G(6)-1/500>47/10.                                \tag{4}
\]
For the last inequality, $\log6>179/100$ suffices.

For $|s|\ge5$ and $|w|\le1$, differentiation of
$\log(1+w/s^6)$ as an absolutely convergent power series shows
$|\Phi_{sss}|<1$. Indeed, putting $x=5^{-6}$, a rational upper bound is
\[
 \frac1{125}\left(12+
 \frac{216x(1+x)}{(1-x)^3}+
 \frac{108x}{(1-x)^2}+\frac{12x}{1-x}\right)<1.          \tag{5}
\]
The entire segment $s_0+[-r,r]$ has modulus greater than five, so
\[
 \Phi(s_0+t,w)=\Phi_0-at^2/2+U(t),\qquad
 |U(t)|\le|t|^3/6 \quad (|t|\le r).
\]
Using $|e^v-1|\le |v|e^{|v|}$ and (3), the local integral differs from
$e^{n\Phi_0}\int_{-r}^r e^{-nat^2/2}\,dt$ by at most
\[
 e^{nS}\frac n6\int_{\mathbb R}|t|^3
       e^{-n(1/14-1/24)t^2}\,dt
 =e^{nS}\frac{4704}{25n}<e^{nS}\frac{189}{n}.           \tag{6}
\]
Completing the Gaussian integral adds an error at most
\[
              e^{nS}\frac{56}{n}e^{-n/224}.            \tag{7}
\]
Since $\operatorname{Re}a>0$, its complete value is
$e^{n\Phi_0}\sqrt{2\pi/(na)}$, whose modulus exceeds
$5e^{nS}/\sqrt n$ by $|a|<1/5$ and $\pi>5/2$.

\section{The complete contour and all its tails}
For each fixed $n,w$, the integrand $(w+s^6)^ne^{-ns}$ is entire.
A rectangle with right side tending to positive infinity therefore gives
\[
 I_n(w)=\int_0^{i\beta}(w+s^6)^ne^{-ns}\,ds
       +\int_0^\infty(w+(x+i\beta)^6)^ne^{-n(x+i\beta)}\,dx. \tag{8}
\]
The right side of the rectangle vanishes because its exponential decay
dominates the fixed polynomial. This applies for either sign of $\beta$;
no branch cut or neglected contour piece is involved. The local integral
above is precisely the part with $|x-\alpha|\le r$.

On the vertical segment, $|\beta|\le\delta$ gives the upper bound
\[
 \delta(1+\delta^6)^n<e^{nS}\delta e^{-4n}.             \tag{9}
\]
On $0\le x\le4$, put $u=x+\delta$. Since $u^6e^{-u}$ is increasing for
$u<6$,
\[
 |w+(x+i\beta)^6|e^{-x}
 \le (1+u^6)e^{-x}
 \le e^\delta+(4+\delta)^6e^{-4}<83<e^{23/5}.
\]
The two numerical comparisons follow from $e^\delta\le1000/999$,
$e>8/3$, and
$(8/3)^4\sum_{j=0}^4(3/5)^j/j!>83$.
Together with (4), this interval contributes at most
\[
                         4e^{nS}e^{-n/10}.             \tag{10}
\]

For $x\ge4$, the pointwise logarithmic estimate is
\[
 \log|w+(x+i\beta)^6|-x
 \le G(x)+\frac{3\delta^2}{x^2}+\frac1{x^6}
 <G(x)+1/4000.                                        \tag{11}
\]
At a zero of its argument the logarithmic upper bound is interpreted by
continuity. For $4\le x\le8$, $G''(x)=-6/x^2\le-3/32$, so
\[
 G(x)\le G(6)-\frac3{64}(x-6)^2.
\]
Outside the local interval, $|x-6|\ge r-\delta>6/25$.
The loss relative to $S$ in (11) is therefore greater than
\[
 \frac3{64}(6/25)^2-\frac1{4000}-\frac1{500}
 =\frac9{20000}>\frac1{2500}.
\]
This part contributes at most $4e^{nS}e^{-n/2500}$.
For $x\ge8$, $G'(x)\le-1/4$ and
$G(6)-G(8)\ge3/16$. Equations (4) and (11) then give the tail bound
\[
                   e^{nS}\frac4n e^{-n/10}.            \tag{12}
\]
Combining (6)--(12) bounds the error in the first part of (2) by $e^{nS}$
times
\[
 \frac{189}{n}+\frac{56}{n}e^{-n/224}
 +4e^{-n/2500}+\frac4n e^{-n/10}+4e^{-n/10}
 +\delta e^{-4n}.                                     \tag{13}
\]
For $n\ge10^6$, multiply (13) by $\sqrt n$ and use
$e^v\ge v^2/2$ and $e^v\ge v$ to bound it by
\[
 \frac{245}{1000}
 +\frac{8\cdot2500^2}{10^9}
 +\frac4{1000}+\frac{800}{10^9}
 +\frac1{4\cdot10^6}<1.
\]
Thus (2) holds uniformly on the closed disk, proving the inner assertion
of Theorem~\ref{thm:annulus}.

\section{Uniform separation of all formal saddle data}
\begin{corollary}\label{cor:saddles}
On $1\le|w|\le2800$, the six roots of $w=s^5(6-s)$ are simple. Each has
$1/2<|s|<7$. One lies in $|s-6|<3/4$ and the other five in $|s|<23/5$.
For distinct saddles $s,t$,
\[
 |s-t|>10^{-13},\qquad
 |\Phi_{ss}(s,w)|>1/4000,\qquad
 \left|\frac{6-s}{6w}-\frac{6-t}{6w}\right|>10^{-18}.  \tag{14}
\]
The endpoint phase derivative satisfies
\[
 \left|\frac1w-\frac{6-s}{6w}\right|>1/40000.           \tag{15}
\]
These are bounds on saddle data, not on distances between polynomial roots.
\end{corollary}
\begin{proof}
The two Rouch\'e circles are certified by
$(21/4)^5(3/4)>2800$ and $(23/5)^5(7/5)>2800$.
The lower bound $|s|>1/2$ follows from
$|s^5(6-s)|\le(1/2)^5(13/2)<1$ when $|s|\le1/2$.
For $H_w(s)=s^6-6s^5+w$,
\[
 \operatorname{Disc}_s H_w=6^6w^4(3125-w),\qquad
 |\operatorname{Disc}_s H_w|\ge6^6\cdot325.
\]
All root differences are less than fourteen, so any particular difference
$d$ satisfies $|d|^2 14^{28}>6^6\cdot325$. Since
$6^6\cdot325\cdot10^{26}>14^{28}$, the separation in (14) follows.
On the segment joining $5$ to $s$, the derivative of $s^5(6-s)$ has
modulus at most $6\cdot7^4\cdot12$. Thus
\[
 |s-5|\ge\frac{325}{6\cdot7^4\cdot12},\qquad
 |\Phi_{ss}(s,w)|=\frac{|s-5|}{|s|}>\frac1{4000}.
\]
Along a saddle branch the total $w$ derivative of its phase equals
$1/(w+s^6)=(6-s)/(6w)$. The remaining bounds follow from
$|w|\le2800$, $|s-t|>10^{-13}$ and $|s|>1/2$.
\end{proof}

\section{The remaining simplicity obligation}
The annulus removes both critical values $0$ and $3125$ from every
possible root region for $n\ge10^6$. It does not determine which saddles
contribute to the original contour, establish a global two-contribution
formula, or control cancellations where several contributions are comparable.
Endpoint contributions at nonzero $w$ remain important even though the
singular endpoint value $w=0$ is now excluded.

Here is an exact sufficient local criterion. Suppose, on a neighborhood,
$I_n=U+V+E$ with analytic nowhere-zero $U,V$. Put
$\ell_U=U'/U$, $\ell_V=V'/V$. At any common zero of $I_n,I_n'$,
\[
             (\ell_U-\ell_V)U=\ell_V E-E'.            \tag{16}
\]
Consequently the strict bound
\[
 |\ell_U-\ell_V|\,|U|>|E'|+|\ell_V|\,|E|             \tag{17}
\]
excludes a multiple zero there. For saddle or endpoint models,
(14)--(15) separate the leading phase derivatives. Effective bounds on
amplitudes and differentiated remainders are still needed to verify (17).
A proof must also cover all root neighborhoods, including any region
requiring more than two comparable contributions, or give a different
uniform $P_n/P_n'$ exclusion. No global simplicity theorem or effective
simplicity cutoff is asserted. The finite prefix below any eventual
simplicity cutoff has not been verified by this note.

The companion exact checker verifies the symbolic identities and rational
constant inequalities used above. It does not formally verify contour
analysis, compute polynomial zeros up to $10^6$, or provide an independent
proof replay. Independent, external, and formal assurance remain false.
\end{document}
